\documentclass[10pt,a4paper]{amsart}
\usepackage[margin=19mm]{geometry}
\usepackage{amsmath,amssymb,mathtools}
\usepackage[scr=rsfs,cal=euler]{mathalfa}
\usepackage{xfrac}
\usepackage[unicode,hidelinks,bookmarks=false]{hyperref}
\newcommand{\ie}{i.e.\ }
\newcommand{\cf}{cf.\ }
\newcommand{\resp}{respectively\ }
\newcommand{\Subsection}{\subsection}
\providecommand{\nobreakdash}{\nobreak\hbox{-}}
\setlength{\emergencystretch}{2em}
\multlinegap0pt
\usepackage{amssymb}
\usepackage{mathtools}
\usepackage{xcolor}
\newtheorem{theorem}{Theorem}
\newcommand{\Dstar}{\mathcal{D}^*}
\newcommand{\ICG}{\mathrm{ICG}}
\newcommand{\vp}{\varphi}
\title{Graph Energy Maximisation for Integral Circulant Graphs of Order $n = p^2q^3$}
\author{Diego Roldan}
\date{Source: arXiv:2604.09491v1 (CC BY 4.0)}
\begin{document}
\maketitle

\noindent\emph{Source and licence.} Graph Energy Maximisation for Integral Circulant Graphs of Order $n = p^2q^3$. Version arXiv:2604.09491v1 (CC BY 4.0), distributed by arXiv under the Creative Commons Attribution 4.0 International licence, http://creativecommons.org/licenses/by/4.0/, as recorded in the arXiv licence metadata for this exact version and shown on the abstract page https://arxiv.org/abs/2604.09491v1. Full licence notice: \texttt{licenses/arxiv-2604.09491-CC-BY-4.0.txt}.\par\smallskip

\noindent\textit{Editorial scope.} Counted unit: the article's theorem thm:kronecker (source0.tex lines 181--198). The article's complete original proof of it is delivered with it (source0.tex lines 363--399, one \texttt{proof} environment of 37 lines, which the article places away from the statement). No mathematics is written, repaired or re-derived: the statement and the proof are the article's own lines, in the article's own notation, and no retained line is substituted or re-flowed.  Retained uncounted as the source-local context the proof needs: lines 207--215; lines 209--214.  Nothing was omitted from any retained span, so the package carries no omission record.  No retained line contains a citation, so the wrapper carries no bibliography and no bibliography entry.  No retained line contains a figure environment, an image-inclusion command or drawing code, so no image is delivered and the package declares no asset dependency.  Editorial formatting only, in addition: an AMS \texttt{amsart} preamble that restates the article's own declarations and macros used by the retained lines, namely the environments \texttt{theorem} on separate counters, together with the standard packages those lines need.  The retained environments print in this document as Theorem 1 in order of appearance, whereas the article prints the same items with the numbering of its own full document.  The complete native source of the article as distributed in the arXiv e-print for this exact version is bundled unchanged as \texttt{sources/198211/source0.tex}.
\begin{theorem}[Closed-form energy]\label{thm:formula}
For every pair of distinct odd primes $(p,q)$,
\begin{align}
  E(p^2q^3,\Dstar)
  &= (5p^2-8p+4)(q-1)(3q^2-2q+1) \notag\\
  &\quad + (p-1)(2p-1)(q^3-2q^2+2q-2) \label{eq:Emax}\\
  &\quad + (p^2-2p+2)(q-1)(q^2+1)+(p-1)q^3. \notag
\end{align}
\end{theorem}

\begin{align}
  E(p^2q^3,\Dstar)
  &= (5p^2-8p+4)(q-1)(3q^2-2q+1) \notag\\
  &\quad + (p-1)(2p-1)(q^3-2q^2+2q-2) \\
  &\quad + (p^2-2p+2)(q-1)(q^2+1)+(p-1)q^3. \notag
\end{align}

\begin{theorem}[Kronecker factorisation]\label{thm:kronecker}
Let $\Dstar=\{1,p^2,pq,q^2,p^2q^2,pq^3\}$ and write $d_{a,b}=p^aq^b$
for $(a,b)\in\{0,1,2\}\times\{0,1,2,3\}$.  Define the partial alternating
geometric sums $S_b(q)=\sum_{k=0}^{b}(-q)^k$, so $S_0=1$, $S_1=1-q$,
$S_2=1-q+q^2$, and set $\alpha_0=-2$, $\alpha_1=2(p-1)$,
$\alpha_2=-(p^2-2p+2)$.  Then for every
$(a,b)\in\{0,1,2\}\times\{0,1,2,3\}$:
\begin{alignat}{2}
  \lambda_{d_{a,b}}(\Dstar) &= \alpha_a\cdot S_b(q),
  &&\quad b=0,1,2,\;\text{all }a,\label{eq:lam-b012}\\
  \lambda_{d_{0,3}}(\Dstar) &= q^3-2q^2+2q-2,
  &&\quad \label{eq:lam-03}\\
  \lambda_{d_{1,3}}(\Dstar) &= (p-1)(2-2q+2q^2-q^3),
  &&\quad \label{eq:lam-13}\\
  \lambda_{d_{2,3}}(\Dstar) &= (p^2-2p+2)(q-1)(q^2+1)+(p-1)q^3.
  &&\quad \label{eq:lam-23}
\end{alignat}
\end{theorem}

\begin{proof}[Proof of Theorem~\ref{thm:formula}]
Given Theorem~\ref{thm:kronecker}, the energy of $\ICG(p^2q^3,\Dstar)$
decomposes as
\[
  E(p^2q^3,\Dstar)
  = \sum_{a=0}^{2}\sum_{b=0}^{3}
    \vp(p^{2-a})\,\vp(q^{3-b})\,
    \bigl|\lambda_{d_{a,b}}(\Dstar)\bigr|.
\]

\textit{Contribution $b\le2$.}  Since $q\ge3$ we have $|S_0|=1$,
$|S_1|=q-1$, $|S_2|=q^2-q+1$.  Computing
\[
  \sum_{b=0}^{2}\vp(q^{3-b})|S_b|
  = q^2(q-1)\cdot1 + q(q-1)(q-1) + (q-1)(q^2-q+1)
  = (q-1)(3q^2-2q+1).
\]
For the $a$-sum, using $|\alpha_a|=(2,\,2(p-1),\,p^2-2p+2)$:
\[
  \sum_{a=0}^{2}\vp(p^{2-a})|\alpha_a|
  = p(p-1)\cdot2 + (p-1)\cdot2(p-1) + 1\cdot(p^2-2p+2)
  = 5p^2-8p+4.
\]
The $b\le2$ contribution is therefore $(5p^2-8p+4)(q-1)(3q^2-2q+1)$.

\textit{Contribution $b=3$.}  Since $q^3-2q^2+2q-2>0$ for $q\ge3$ and
$\vp(q^0)=1$:
\begin{align*}
  &\vp(p^2)|\lambda_{d_{0,3}}|+\vp(p^1)|\lambda_{d_{1,3}}|
   +\vp(p^0)|\lambda_{d_{2,3}}|\\
  &= p(p-1)(q^3-2q^2+2q-2)+(p-1)^2(q^3-2q^2+2q-2)
   +(p^2-2p+2)(q-1)(q^2+1)+(p-1)q^3\\
  &= (p-1)(2p-1)(q^3-2q^2+2q-2)+(p^2-2p+2)(q-1)(q^2+1)+(p-1)q^3.
\end{align*}

Adding the two contributions gives~\eqref{eq:Emax}.\qed
\end{proof}

\end{document}
